\section{Responsibility Division across Layers}

Each layer of the aforementioned Clean architecture covers a separate part of the project, with its own responsibilities.
The layers and their respective responsibilities are as follows:

\begin{itemize}
\item \textbf{Domain}: The domain of the project is tutorial scheduling, and the domain layer is responsible for defining the entities, data value objects and aggregates that the scheduler consists of. Mainly this includes Students, TAs, Lecturers, Modules, Rooms and the Scheduler;
\item \textbf{Use Cases}: The use cases layer is responsible for behaviors of objects in the domain layer. This includes the actual tutorial scheduling itself, user registration, users filling in or updating schedules, as well as creating and updating module and room details;
\item \textbf{Adapters/Interfaces}: The interfaces layer is responsible for the communication between the domain and the infrastructure, which is done through a repository system. This layer also includes APIs that the Web UI uses to communicate with the backend.
\item \textbf{UI + Infrastructure}: This layer contains the infrastructure used, mainly an SQL database for persistent data storage, as well as the user-facing Web UI consisting of a React.js-based frontend.
\end{itemize}
